\section{Working out what each run came to}

Preprocessing hands on a list of runs the lab can answer for. This stage turns each one into a
single row of numbers. The class that drives it is \texttt{Calculator}, in
\path{src/limodel/analyzer/calculate/calculator.py}, and its one public method is
\texttt{find\_results(runs)}. Every row it produces is a \texttt{Result}: a frozen dataclass
holding the run code, the membrane, the two settings the run was done at, the two outcomes the
model is fitted to, and the five further fields the report calls diagnostics. One of those five,
the highest voltage, is fitted as well, because a gate is placed on it.

Nothing in a \texttt{Result} is copied out of a CSV file. Every field is worked out from tank
weights, lab concentrations and the power log. The chain below is given in dependency order,
because that is the order the code reads in: the later quantities call the earlier ones rather
than re-deriving them.

Each quantity is one small class with static methods and no state. They are in
\path{src/limodel/analyzer/calculate/quantities/}, one file per class, and the file name is the
class name. The two constants they all lean on sit in \texttt{constants.py}:
$F = 96485\unit{C/mol}$, $M_\mathrm{Li} = 6.94\unit{g/mol}$,
$M_\mathrm{Mg} = 24.305\unit{g/mol}$, and a solution density of $1000\unit{g/L}$.

\subsection{Volume in a tank}

\textbf{\texttt{VolumeInATank}} --- \texttt{find\_litres(reading, tare\_g)} $\rightarrow$
litres.

Nothing on the rig measured volume. Each tank sat on a balance for the whole run and was weighed
empty beforehand, so the liquid is the balance reading less that empty weight:

\begin{equation}
  V(t) \;=\; \frac{w(t) - w_\mathrm{tare}}{\rho},
  \qquad \rho = 1000\unit{g/L}
  \label{eq:volume}
\end{equation}

$w(t)$ is \texttt{weight\_g} on the \texttt{TankReading} for that minute, and $w_\mathrm{tare}$
comes from the run's own \texttt{tare\_conc\_g}, \texttt{tare\_dil\_g} or \texttt{tare\_rinse\_g}
column, picked by \texttt{TARE\_FIELD\_BY\_TANK}. Taking $\rho$ as exactly
$1000\unit{g/L}$ is a deliberate simplification: these are dilute salt solutions, so a litre
weighs a kilogram closely enough that grammes on the balance can be read straight off as litres.
Returns \texttt{None} when the balance gave nothing for that minute, or the tare column is empty.

\subsection{Amount of a metal in a tank}

\textbf{\texttt{AmountInATank}} --- \texttt{find\_milligrams(reading, tare\_g, metal)}
$\rightarrow$ mg.

The lab reports a concentration, which on its own says nothing useful here. Every tank gains and
loses liquid through the run, so a concentration can rise while the metal in the tank falls. Only
the mass can be compared minute to minute:

\begin{equation}
  m_X(t) \;=\; c_X(t)\,V(t),
  \qquad [\unit{mg/L}] \times [\unit{L}] = [\unit{mg}]
  \label{eq:amount}
\end{equation}

$c_X(t)$ is read off the \texttt{Sample} hanging on that reading, through the \texttt{Metal}'s
own \texttt{field\_name} --- \texttt{li\_mg\_l} for lithium, \texttt{mg\_mg\_l} for magnesium.
Returns \texttt{None} when the minute has no sample at all, when the lab left that metal's cell
blank, or when \eqref{eq:volume} could not be worked out.

\subsection{Amount that crossed}

\textbf{\texttt{AmountThatCrossed}} --- \texttt{find\_milligrams(run, tank, metal)}
$\rightarrow$ mg.

How much of a metal a tank ended with, less how much it started with:

\begin{equation}
  \Delta m_X^{\,\mathrm{tank}} \;=\; m_X(t_\mathrm{last}) - m_X(t_\mathrm{first})
  \label{eq:crossed}
\end{equation}

Only sampled minutes can take part, since \eqref{eq:amount} needs a lab result. The ends are
found by \emph{minute}, not by position in the list: \texttt{find\_amounts\_by\_minute} builds a
dictionary keyed by minute, keeps only the minutes that produced a number, and the ends are taken
from \texttt{sorted()} of those keys. That matters because readings arrive in file order and
because preprocessing may have removed a minute in the middle --- or at an end, as it did on
AJ's concentrate at minute~0.

\texttt{find\_ends} returns the pair $(m_X(t_\mathrm{first}), m_X(t_\mathrm{last}))$ rather than
the difference, because recovery, the feed ratio and the balance gap all need the \emph{start}
on its own. Both methods return \texttt{None} when fewer than two sampled minutes of that tank
came through with a number, and \texttt{find\_ends} also returns \texttt{None} for a tank name
it does not know.

\subsection{Lithium recovery}

\textbf{\texttt{LithiumRecovery}} --- \texttt{find\_percent(run)} $\rightarrow$ per cent.

The share of the lithium the feed started with that ended up in the concentrate:

\begin{equation}
  R_\mathrm{Li} \;=\; 100 \times
  \frac{\Delta m_\mathrm{Li}^{\,\mathrm{conc}}}{m_\mathrm{Li}^{\,\mathrm{dil}}(t_\mathrm{first})}
  \label{eq:recovery}
\end{equation}

Milligrams stand on both the top and the bottom, so $M_\mathrm{Li}$ cancels. There is nothing to
gain by converting to moles, and one fewer place for a unit slip to hide. The result is a per
cent, not a fraction, because that is how the gate is written and how the report prints it.

The bottom is the dilute's \emph{first} sampled minute, not what it ended with. What the feed
held before anything moved is the whole prize; measuring against the end would make a run that
moved nothing look like a run that moved everything.

The top is a difference, \eqref{eq:crossed}, and that subtraction is not cosmetic. The
concentrate never starts with zero lithium, and across these runs its starting amount spans a
factor of more than five hundred, depending on which salt the concentrate was made up with.
Over the sixteen runs the calculator uses it runs from $2.6\unit{mg}$ up to $1455.1\unit{mg}$.
The low end is AN and AO, whose concentrate was made up on a sodium salt rather than a lithium
one, so only traces were there at the start ($18.6$ and $2.6\unit{mg}$). The high end is AP and
AQ, started on LiCl at $5\unit{g/L}$ ($1414.4$ and $1455.1\unit{mg}$). The twelve runs started
on LiCl at $0.5\unit{g/L}$ sit between $126.4$ and $161.7\unit{mg}$. On AQ the lithium already
in the
concentrate is about double what the whole feed held, so treating the start as zero would not
merely bias recovery --- it would report roughly $214\,\%$ in place of $11.4\,\%$.

As a check on \eqref{eq:recovery}, run AE: the concentrate gained
$309.55\unit{mg}$ of lithium, the dilute began with $808.74\unit{mg}$, and the report prints
$38.3\,\%$.

\subsection{Charge that passed}

\textbf{\texttt{ChargeThatPassed}} --- \texttt{find\_coulombs(run)} $\rightarrow$ C, and
\texttt{find\_mah\_per\_cm2(run)} $\rightarrow$ mA\,h/cm\textsuperscript{2}.

No coulomb counter was fitted. The supply current was written down every so many minutes, so the
charge is the area under those readings, by the trapezium rule over consecutive logged minutes:

\begin{equation}
  Q \;=\; 60 \sum_{k} \frac{I_k + I_{k+1}}{2}\,\bigl(t_{k+1} - t_k\bigr)
  \label{eq:charge}
\end{equation}

The $60$ converts amp-minutes to amp-seconds, and an amp-second is a coulomb. Spread over the
membrane area and converted to the unit the field quotes:

\begin{equation}
  q \;=\; \frac{Q}{3600\,A} \times 1000
  \qquad [\unit{mA\,h/cm^2}]
  \label{eq:charge-density}
\end{equation}

The gaps $t_{k+1}-t_k$ are whatever the log holds; nothing is resampled onto an even grid, and no
reading is carried forward past the last logged minute. Only minutes with a
\texttt{supply\_current\_a} take part, so a minute preprocessing removed simply widens the
trapezium either side of it. \texttt{find\_coulombs} returns \texttt{None} on fewer than two such
minutes, and \texttt{find\_mah\_per\_cm2} also returns \texttt{None} when
\texttt{area\_cm2} is zero or missing. AE gives $Q = 1178.55\unit{C}$ over
$A = 6\unit{cm^2}$, which is $54.6\unit{mA\,h/cm^2}$.

\subsection{Current density}

\textbf{\texttt{CurrentDensity}} --- \texttt{find\_ma\_per\_cm2(run)} $\rightarrow$
mA/cm\textsuperscript{2}.

\begin{equation}
  J \;=\; \frac{1000 \times \operatorname{median}\bigl(\{I_k\}\bigr)}{A}
  \qquad [\unit{mA/cm^2}]
  \label{eq:current-density}
\end{equation}

This is the first of the two columns the polynomial is fitted over, so how it is obtained
matters more than it looks. It comes from the \emph{measured} supply current, not from the set
point: the \texttt{set\_current\_ma} column is empty on all twenty-one runs in \path{data/runs.csv},
so there is no set point to read. The median rather than the mean, because a run was held at one
current throughout and the number that labels the whole run should not be dragged by a single
misread minute. Returns \texttt{None} when no minute carries a current, or the run has no area.

\subsection{Stack voltage}

\textbf{\texttt{StackVoltage}} --- \texttt{find\_highest\_volts(run)} $\rightarrow$ V, and
\texttt{find\_starting\_volts(run)} $\rightarrow$ V.

\begin{equation}
  U_\mathrm{max} \;=\; \max_k U_k,
  \qquad
  U_\mathrm{start} \;=\; U_{k_\mathrm{min}}
  \label{eq:voltage}
\end{equation}

$U_\mathrm{max}$ is the highest \texttt{supply\_voltage\_v} the run reached, and it is the figure
the $20\unit{V}$ gate is applied to: a ceiling is only cleared if it is cleared at the worst
minute, not on average. $U_\mathrm{start}$ is the voltage at the earliest minute that has one ---
the earliest, deliberately, not the lowest --- and says what the stack asked for before anything
had crossed. Only $U_\mathrm{max}$ reaches a \texttt{Result}; $U_\mathrm{start}$ is there for
reading a run. Both return \texttt{None} when the power log never recorded a voltage.

\subsection{Transport numbers}

\textbf{\texttt{TransportNumbers}} --- \texttt{find\_numbers(run)} $\rightarrow$
$(t_\mathrm{Li}, t_\mathrm{Mg})$, dimensionless, summing to $1$.

A transport number is a share of the \emph{charge} carried, not a share of the ions. So each
metal's crossing is turned into millimoles and then weighted by the charge its own ion carries:
$z_\mathrm{Li} = 1$, $z_\mathrm{Mg} = 2$. Those charges live on the \texttt{Metal} dataclass
alongside the molar mass, which is why they travel together.

\begin{equation}
  n^\star_X \;=\; \frac{z_X\,\Delta m_X^{\,\mathrm{conc}}}{M_X},
  \qquad
  t_\mathrm{Li} = \frac{n^\star_\mathrm{Li}}{n^\star_\mathrm{Li} + n^\star_\mathrm{Mg}},
  \qquad
  t_\mathrm{Mg} = \frac{n^\star_\mathrm{Mg}}{n^\star_\mathrm{Li} + n^\star_\mathrm{Mg}}
  \label{eq:transport}
\end{equation}

Why magnesium counts twice: the current through the stack is carried by charge, and one
Mg\textsuperscript{2+} crossing the membrane moves two units of charge while one
Li\textsuperscript{+} moves one. If the $z_X$ were dropped, these would be shares of the ion
count, and a membrane that let through equal moles of each would look even-handed when in fact
magnesium was doing two-thirds of the work. Returns \texttt{None} when either metal's crossing is
unknown, or when the two weighted amounts sum to zero. On AE,
$(t_\mathrm{Li}, t_\mathrm{Mg}) = (0.680, 0.320)$.

\subsection{The two feed ratios}

\textbf{\texttt{FeedRatio}} has two methods, both returning mol/mol:
\texttt{find\_li\_over\_mg\_in\_the\_dilute(run)} and
\texttt{find\_mg\_over\_li\_from\_the\_setup(run)}.

There are two, kept apart by name because they answer different questions and are not
reciprocals of one another.

The \emph{measured} one is taken from the lab, at the first dilute minute for which both metals
were reported:

\begin{equation}
  r_\mathrm{dil} \;=\;
  \frac{m_\mathrm{Li}^{\,\mathrm{dil}}(t_0)\,/\,M_\mathrm{Li}}
       {m_\mathrm{Mg}^{\,\mathrm{dil}}(t_0)\,/\,M_\mathrm{Mg}}
  \label{eq:feed-measured}
\end{equation}

$t_0$ is the earliest minute present in both metals' dictionaries --- the intersection of the
two, so a minute where the lab reported only one metal is skipped rather than half-used. Both
amounts come from the same tank at the same minute, so the litres in \eqref{eq:amount} cancel and
the milligrams carry the ratio on their own. This is the one permselectivity divides by.

The \emph{set up} one is taken from the run's own feed columns, and is the other way up:

\begin{equation}
  r_\mathrm{setup} \;=\;
  \frac{c_\mathrm{Mg}^{\,\mathrm{feed}}\,/\,M_\mathrm{Mg}}
       {c_\mathrm{Li}^{\,\mathrm{feed}}\,/\,M_\mathrm{Li}}
  \label{eq:feed-setup}
\end{equation}

using \texttt{mg\_feed\_mg\_l} and \texttt{li\_feed\_mg\_l}. This is the second of the two
columns the polynomial is fitted over, and it is the intended ratio rather than the measured one
on purpose: it is a dial setting that was chosen before the run, known exactly, and the same for
every replicate at that point. Using a measured value would put lab noise into a design column
and spread what should be one design point into a cloud. It goes into \texttt{Result} as
\texttt{feed\_mg\_li\_ratio}, which is why the report's column is headed Mg/Li. Each returns
\texttt{None} when its inputs are missing or its denominator is zero. On AE,
$r_\mathrm{dil} = 0.852$ and $r_\mathrm{setup} = 1.428$.

\subsection{Permselectivity}

\textbf{\texttt{Permselectivity}} --- \texttt{find\_value(run)} $\rightarrow$
dimensionless.

The outcome the whole project maximises:

\begin{equation}
  P \;=\; \frac{t_\mathrm{Li} / t_\mathrm{Mg}}{r_\mathrm{dil}}
  \label{eq:perm}
\end{equation}

The top is what the membrane delivered. The bottom is what it was offered. Dividing by
$r_\mathrm{dil}$ is what makes the number mean anything across runs, and the arithmetic is worth
following. Writing $n_X$ for the moles of each metal that crossed, \eqref{eq:transport} gives

\begin{equation}
  \frac{t_\mathrm{Li}}{t_\mathrm{Mg}} \;=\;
  \frac{z_\mathrm{Li}\,n_\mathrm{Li}}{z_\mathrm{Mg}\,n_\mathrm{Mg}} \;=\;
  \frac{n_\mathrm{Li}}{2\,n_\mathrm{Mg}}
  \label{eq:transport-ratio}
\end{equation}

Now take a membrane with no preference at all, letting the two metals through in proportion to
what the feed holds: $n_\mathrm{Li}/n_\mathrm{Mg} = r_\mathrm{dil}$. Then
\eqref{eq:transport-ratio} is $r_\mathrm{dil}/2$, and \eqref{eq:perm} comes to exactly $0.5$ ---
whatever the feed was. That is the number to read $P$ against: above $0.5$ the membrane preferred
lithium, below it magnesium, and the factor of two is the charge weighting, not a choice.

Without the division the baseline would slide. The same indifferent membrane would score
$r_\mathrm{dil}/2$: that is $0.426$ at AE's feed, and $0.043$ at a feed ten times richer in
magnesium. Two runs at different feeds could not then be set side by side. Because the feed ratio
is one of the two model inputs, this is not a detail. The polynomial is asked how $P$ changes as
the feed ratio moves; if the baseline moved with it too, the fitted feed-ratio term would be
measuring the normalisation rather than the membrane, and every conclusion about feed sensitivity
would be an artefact.

Returns \texttt{None} when the transport numbers or the measured feed ratio are missing, or when
$t_\mathrm{Mg}$ or $r_\mathrm{dil}$ is zero. On AE,
$P = (0.680/0.320)/0.852 = 2.498$, about five times indifference.

\subsection{Energy per mole of lithium}

\textbf{\texttt{EnergyPerMoleOfLithium}} --- \texttt{find\_joules(run)} $\rightarrow$ J, and
\texttt{find\_watt\_hours(run)} $\rightarrow$ W\,h/mol.

Power is multiplied minute by minute before any averaging, because voltage and current do not
move together and the average of a product is not the product of the averages. The same
trapezium rule as \eqref{eq:charge} then runs over the minutes that carry both:

\begin{equation}
  E \;=\; 60 \sum_{k} \frac{U_k I_k + U_{k+1} I_{k+1}}{2}\,\bigl(t_{k+1} - t_k\bigr)
  \qquad [\unit{J}]
  \label{eq:energy}
\end{equation}

\begin{equation}
  n_\mathrm{Li} \;=\; \frac{\Delta m_\mathrm{Li}^{\,\mathrm{conc}}}{1000\,M_\mathrm{Li}},
  \qquad
  e \;=\; \frac{E}{3600\,n_\mathrm{Li}}
  \qquad [\unit{W\,h/mol}]
  \label{eq:energy-per-mole}
\end{equation}

Per mole of lithium delivered, rather than per hour or per litre, so a short run and a long one,
a gentle current and a hard one, can be compared on the thing being paid for. Note the set of
minutes here is narrower than in \eqref{eq:charge}: a minute with a current but no voltage counts
towards charge and not towards energy. \texttt{find\_joules} returns \texttt{None} on fewer than
two minutes carrying both. \texttt{find\_watt\_hours} also returns \texttt{None} when
$n_\mathrm{Li} \le 0$, since a concentrate that gained no lithium, or lost some, has no per-mole
cost to quote. AE comes to $E = 14759.8\unit{J}$ and $e = 91.9\unit{W\,h/mol}$.

\subsection{The balance gap}

\textbf{\texttt{BalanceGap}} --- \texttt{find\_percent(run, metal)} $\rightarrow$ per cent,
computed once for lithium and once for magnesium.

\begin{equation}
  L_X \;=\; m_X^{\,\mathrm{dil}}(t_\mathrm{first}) - m_X^{\,\mathrm{dil}}(t_\mathrm{last}),
  \qquad
  G_X \;=\; 100 \times
  \frac{L_X - \bigl(\Delta m_X^{\,\mathrm{conc}} + \Delta m_X^{\,\mathrm{rinse}}\bigr)}{L_X}
  \label{eq:gap}
\end{equation}

This corrects nothing. No other quantity reads it, no fit uses it, and no run is dropped on it.
It exists so a reader can judge how much weight one run's numbers deserve.

A large positive gap means the dilute lost metal that was not found in either of the other two
tanks --- the readings do not add up, and the run's recovery and permselectivity should be read
with suspicion. A negative gap means more was found than the dilute is recorded as having lost,
which usually points at a weight or a concentration being wrong at one end. AT shows
$-272.5\,\%$: its dilute barely changed, losing $5.05\unit{mg}$, while its concentrate gained
$18.82\unit{mg}$, so a small absolute mismatch divides by a tiny denominator. AT's recovery of
$2.3\,\%$ should be read in that light.

The liquid carried off in the sample bottles is \emph{not} accounted for, by choice. Every sample
drawn removes both volume and metal from a tank for good. Rather than model that removal --- which
would mean tracking bottle volumes nobody recorded reliably --- it is left out, and the gap
absorbs it. So a modest positive gap is expected, and the gap is a bound on the error rather than
a measurement of it. A missing rinse figure is likewise counted as zero gained, which can only
make the gap look worse; for a diagnostic, erring towards alarm is the honest direction. Fifteen
of the sixteen used runs do have a rinse figure. Returns \texttt{None} when the dilute ends or
the concentrate's gain cannot be worked out, or when $L_X \le 0$: a share of a loss means nothing
where there was no loss, and a dilute that gained would turn the sign of the gap over.

\subsection{Missing is never a guess}

Every method above returns \texttt{None} when it cannot be worked out. Nothing is filled in with
a zero, a mean, or a carried-forward value. A zero is a claim about the chemistry and would be
fitted as one; \texttt{None} is a claim about the records and can be counted.

\texttt{Calculator.\_find\_result} then applies one rule. Permselectivity and lithium recovery
are what the model is fitted to, so a run missing either has nothing to contribute: the drop is
written in the journal with the reason naming which of the two failed, and \texttt{None} comes
back instead of a \texttt{Result}. Everything else may be \texttt{None} on a run that is
otherwise sound, and that is not a reason to let the run go --- the report prints a dash in the
cell and names the run under \emph{diagnostics missing} at the end. On the current data all
sixteen surviving runs produce a complete row, so no drop happens at this stage and no dash is
printed; the five runs the report says were dropped were dropped during preprocessing.

\subsection{Every quantity in one table}

\begin{table}[htbp]
  \centering
  \caption{The quantities, in dependency order. Only the fields of \texttt{Result} reach the
    model or the report.}
  \label{tab:quantities}
  \small
  \begin{tabular}{@{}l l l p{0.30\linewidth}@{}}
    \toprule
    Quantity & Class & Unit & Used for \\
    \midrule
    Volume in a tank & \texttt{VolumeInATank} & $\unit{L}$
      & turning a concentration into a mass \\
    Amount in a tank & \texttt{AmountInATank} & $\unit{mg}$
      & one tank at one minute \\
    Amount that crossed & \texttt{AmountThatCrossed} & $\unit{mg}$
      & recovery, transport numbers, feed ratio, energy, balance gap \\
    Lithium recovery & \texttt{LithiumRecovery} & \%
      & fitted outcome; the recovery gate \\
    Charge that passed & \texttt{ChargeThatPassed} & $\unit{mA\,h/cm^2}$
      & diagnostic: how hard the run was driven in total \\
    Current density & \texttt{CurrentDensity} & $\unit{mA/cm^2}$
      & model input $J$; the setting the search optimises \\
    Stack voltage & \texttt{StackVoltage} & $\unit{V}$
      & fitted outcome; the $20\unit{V}$ gate \\
    Transport numbers & \texttt{TransportNumbers} & ---
      & permselectivity \\
    Feed ratio, measured & \texttt{FeedRatio} & $\unit{mol/mol}$
      & the denominator of permselectivity \\
    Feed ratio, set up & \texttt{FeedRatio} & $\unit{mol/mol}$
      & model input $R$; the second search setting \\
    Permselectivity & \texttt{Permselectivity} & ---
      & fitted outcome; what the search maximises \\
    Energy per mole Li & \texttt{EnergyPerMoleOfLithium} & $\unit{W\,h/mol}$
      & diagnostic: the electricity cost of a run \\
    Balance gap, Li and Mg & \texttt{BalanceGap} & \%
      & diagnostic only; corrects nothing \\
    \bottomrule
  \end{tabular}
\end{table}
